\chapter{Implementation}
\label{ch:implementation}

This chapter presents the software implementation of the Multi-Robot Calibration (MRC) system.
First, the overall system architecture and its integration into ROS\,2 are described
(\autoref{sec:system_architecture}).
Subsequently, the individual software components are presented in detail:
the data acquisition pipeline (\autoref{sec:data_acquisition}),
the point cloud registration backends (\autoref{sec:registration_backends}),
the optional ICP refinement stage (\autoref{sec:icp_refinement}),
and the pose graph optimization module (\autoref{sec:pose_graph}).
Finally, \autoref{sec:ros2_interfaces} describes the custom ROS\,2 message, service, and action
interfaces that define the communication protocol between the distributed nodes.

%% ---------------------------------------------------------------------------
\section{System Architecture}
\label{sec:system_architecture}

The MRC system follows a distributed master--slave architecture realized as a set of
ROS\,2 nodes.
\autoref{fig:architecture} provides an overview of the data flow between components.

\begin{itemize}
    \item \textbf{Master node} (\texttt{mrc\_master}):
          Coordinates the calibration workflow, hosts the registration backends,
          optionally runs pose graph optimization, and publishes the resulting
          transforms via the \texttt{tf\_static} broadcaster.
    \item \textbf{Slave nodes} (\texttt{mrc\_slave}):
          Each robot in the fleet runs a slave node that subscribes to the local
          LiDAR sensor topic(s), accumulates and preprocesses point clouds, extracts
          geometric features, and provides the processed data to the master upon
          request through a ROS\,2 service.
\end{itemize}

% TODO: insert architecture diagram
% \begin{figure}[htbp]
%   \centering
%   \includegraphics[width=\textwidth]{figures/mrc_architecture.pdf}
%   \caption{System architecture of the MRC framework showing the master--slave
%            communication pattern and data flow.}
%   \label{fig:architecture}
% \end{figure}

The calibration is triggered either through a \texttt{std\_srvs/Trigger} service
or, preferably, through a ROS\,2 action server that provides asynchronous execution
with progress feedback.
Upon receiving a calibration request, the master sequentially queries each discovered
slave for its registration data, performs pairwise point cloud registration, optionally
refines and optimizes the results, and finally broadcasts the estimated transforms as
static TF frames.

A \emph{factory pattern} is employed for the registration backends, allowing the
algorithm to be selected at runtime via a ROS\,2 parameter without modifying the
source code.
Currently, two backends are implemented: a feature-based method using
KISS-Matcher~\cite{KISSMATCHER} and a dense method using
Fast Robust ICP (FRICP)~\cite{zhang_fast}.

\subsection{Robot Discovery}

The master node discovers robots in the fleet by subscribing to
\texttt{hector\_multi\_robot\_msgs/RobotAnnouncement}\footnote{https://github.com/tu-darmstadt-ros-pkg/hector_multi_robot} messages.
Each slave periodically broadcasts its robot identifier, type, and namespace.
The master maintains a registry of known robots and constructs the
corresponding service client endpoints for data retrieval.

%% ---------------------------------------------------------------------------
\section{Data Acquisition and Preprocessing}
\label{sec:data_acquisition}

The slave node is responsible for acquiring LiDAR data and preparing it for
registration.
It supports two operational modes to accommodate different sensor types:

\begin{description}
    \item[Latest-scan mode (\texttt{accumulation\_type = "none"})]
        Uses only the most recent point cloud message.
        This mode is suitable for repetitive LiDARs.
    \item[Time-window accumulation (\texttt{accumulation\_type = "time\_window"})]
        Accumulates point cloud messages over a configurable time window
        $\Delta t_\mathrm{acc}$ (default: \SI{1.0}{\second}).
        This mode is designed for non-repetitive scanning LiDARs
        (e.g., Livox Mid-360) where a single message covers only a partial
        field of view, and temporal accumulation is required to obtain
        sufficient spatial coverage.
\end{description}

When multiple LiDAR topics are configured per robot, the slave merges the
individual point clouds into a common coordinate frame using TF lookups,
producing a single unified cloud for registration.

\subsection{Feature Extraction}
\label{sec:feature_extraction}

When the KISS-Matcher ~\cite{KISSMATCHER} backend is selected, the slave additionally extracts
geometric features from the accumulated point cloud.
The extraction pipeline consists of the following steps:

\begin{enumerate}
    \item \textbf{Voxel grid downsampling} with a configurable leaf size
          $v_\mathrm{kiss}$ (default: \SI{0.3}{\metre}) to reduce the number
          of points while preserving the geometric structure.
    \item \textbf{Normal estimation} using a neighborhood radius of
          $r_n = \gamma_n \cdot v_\mathrm{kiss}$, where $\gamma_n$ is the
          normal radius gain (default: $\gamma_n = 3.0$).
    \item \textbf{FasterPFH descriptor computation}~\cite{KISSMATCHER} with a
          search radius $r_f = \gamma_f \cdot v_\mathrm{kiss}$
          ($\gamma_f = 5.0$ by default), yielding a 33-dimensional descriptor
          vector for each keypoint.
\end{enumerate}

The computed features (keypoints and descriptors) are serialized into a custom
\texttt{KissFeatures} message and cached on the slave to avoid redundant
recomputation across multiple calibration requests.

%% ---------------------------------------------------------------------------
\section{Registration Backends}
\label{sec:registration_backends}

Both backends implement a common abstract interface \texttt{MRCRegistration}
with a single method \texttt{registerSourceToTarget()}, which takes source and
target data as input and returns a \texttt{RegistrationResult} containing the
estimated rigid body transformation
$\mathbf{T} \in \mathrm{SE}(3)$,
quality metrics (RMSE, inlier count), and a success flag.

\subsection{KISS-Matcher Backend}
\label{sec:kiss_matcher}

The KISS-Matcher ~\cite{KISSMATCHER} backend performs feature-based global registration without
requiring an initial pose estimate.
The algorithm proceeds as follows:

\begin{enumerate}
    \item \textbf{Correspondence establishment:}
          The ROBIN matching module\footnote{https://github.com/MIT-SPARK/ROBIN} matches source and target
          descriptors using a ratio test.
          A configurable maximum number of correspondences
          $N_\mathrm{max}$ and a tuple scale parameter $\tau$ control the
          matching behavior.

    \item \textbf{Robust pose estimation:}
          The matched correspondences are passed to a GNC-TLS
          (Graduated Non-Convexity with Truncated Least Squares) solver
          that jointly estimates the rotation $\mathbf{R} \in \mathrm{SO}(3)$
          and translation $\mathbf{t} \in \mathbb{R}^3$ while being robust
          to a high fraction of outlier correspondences.

    \item \textbf{Quality assessment:}
          The inlier set is determined by the solver.
          The registration is accepted only if the inlier ratio
          $r_\mathrm{inlier} = n_\mathrm{inlier} / n_\mathrm{corr}$
          exceeds a minimum threshold (default: 5\%).
          The translation RMSE is computed on the inlier correspondences
          after applying the estimated transform:
          \begin{equation}
              \mathrm{RMSE}_t = \sqrt{
                  \frac{1}{n_\mathrm{inlier}}
                  \sum_{i \in \mathcal{I}}
                  \lVert \mathbf{R}\,\mathbf{p}_i^s + \mathbf{t} - \mathbf{p}_i^t \rVert^2
              }
          \end{equation}
          where $\mathcal{I}$ denotes the inlier set and $\mathbf{p}_i^s$,
          $\mathbf{p}_i^t$ are corresponding source and target keypoints.
          An approximate rotation RMSE is derived as
          $\mathrm{RMSE}_r \approx \mathrm{RMSE}_t \,/\, \bar{d}$,
          where $\bar{d}$ is the mean lever arm (distance of target inlier
          points from their centroid).
\end{enumerate}

\subsection{FRICP Backend}
\label{sec:fricp}

The FRICP backend performs dense point-to-point registration using the
Fast Robust ICP algorithm~\cite{zhang_fast}.
Unlike KISS-Matcher ~\cite{KISSMATCHER}, it operates directly on the raw point clouds without
feature extraction, but requires a reasonable initial alignment or operates
from the identity transform.

\begin{enumerate}
    \item \textbf{Preprocessing:}
          Both source and target point clouds are converted to Eigen matrices.
          The clouds are demeaned (centered at the origin) as required by the
          FRICP implementation.

    \item \textbf{Registration:}
          The FRICP solver is configured with a Welsch cost function
          (\texttt{ICP::WELSCH}) and Anderson acceleration for faster convergence.
          This robust cost function down-weights large residuals, providing
          tolerance against partial overlap and outliers.

    \item \textbf{Quality assessment:}
          After registration, the nearest-neighbor RMSE is computed on the
          full-resolution point clouds.
          A $3\times$-median outlier rejection scheme is applied:
          for each transformed source point, the nearest-neighbor distance to
          the target cloud is computed; points with distances exceeding three
          times the median distance are classified as outliers and excluded
          from the RMSE computation.
\end{enumerate}

%% ---------------------------------------------------------------------------
\section{ICP Refinement}
\label{sec:icp_refinement}

An optional ICP refinement stage can be applied after the initial registration
from either backend.
This stage uses the PCL\footnote{https://pointclouds.org/} implementation of the standard Iterative Closest Point
algorithm, initialized with the coarse transform as the starting guess.

The refinement proceeds as follows:
\begin{enumerate}
    \item Both point clouds are downsampled using a voxel grid filter to
          reduce computational cost during alignment.
    \item ICP is executed on the downsampled clouds, starting from the coarse
          transform.
    \item The refined transform is evaluated on the \emph{full-resolution}
          clouds using the same nearest-neighbor RMSE metric with
          $3\times$-median outlier rejection as the backends.
    \item The refined result is accepted only if the RMSE improves compared
          to the coarse registration; otherwise, the original result is retained.
\end{enumerate}

This conservative acceptance criterion ensures that the refinement stage can
never degrade the calibration quality.

%% ---------------------------------------------------------------------------
\section{Pose Graph Optimization}
\label{sec:pose_graph}

When calibrating more than two robots, the individual pairwise registrations
can be jointly optimized using pose graph optimization to enforce global
consistency.
The MRC system implements this using the GTSAM
library\footnote{https://github.com/borglab/gtsam}.

\subsection{Graph Construction}

The pose graph models each robot's pose as a variable
$\mathbf{X}_i \in \mathrm{SE}(3)$ and each pairwise registration as a
constraint (factor) between two variables.
Two graph topologies are supported:

\begin{description}
    \item[Star topology:]
        The master robot is anchored at the origin with a tight prior
        ($\sigma_t = 10^{-3}\,\mathrm{m}$, $\sigma_r = 10^{-3}\,\mathrm{rad}$).
        Edges connect only the master to each slave, resulting in
        $n - 1$ registration edges for $n$ robots.
    \item[Mesh topology:]
        In addition to the master--slave edges, all pairwise slave--slave
        registrations are computed and added as edges, yielding up to
        $\binom{n}{2}$ constraints.
        This provides additional redundancy and can improve robustness.
\end{description}

\subsection{Noise Model}
\label{sec:noise_model}

Each edge in the pose graph is assigned a diagonal Gaussian noise model in
the $\mathrm{SE}(3)$ tangent space with six components
$(\sigma_{r_x}, \sigma_{r_y}, \sigma_{r_z}, \sigma_{t_x}, \sigma_{t_y},
\sigma_{t_z})$, following the GTSAM \texttt{Pose3} convention where the
first three components represent rotational uncertainty and the last three
represent translational uncertainty.
The isotropic assumption
$\sigma_{r_x} = \sigma_{r_y} = \sigma_{r_z} \eqqcolon \sigma_r$ and
$\sigma_{t_x} = \sigma_{t_y} = \sigma_{t_z} \eqqcolon \sigma_t$ reduces
the parameterization to two scalar values per edge.

\paragraph{Sigma derivation.}
The sigmas are computed in a three-stage pipeline implemented in
\texttt{deriveSigmas()}:

\begin{enumerate}
    \item \textbf{Initialization.}
          The translational and rotational sigmas are initialized with
          configurable default values
          ($\sigma_t^{(0)} = \SI{0.3}{\metre}$,
           $\sigma_r^{(0)} = \SI{0.2}{\radian}$).
    \item \textbf{RMSE override.}
          If the registration backend provides RMSE estimates---as
          KISS-Matcher does via its GNC-TLS residuals and FRICP via its
          Welsch cost---these values replace the defaults:
          \begin{equation}
              \sigma_t = \begin{cases}
                  \mathrm{RMSE}_t & \text{if available} \\
                  \sigma_t^{(0)}  & \text{otherwise}
              \end{cases}, \qquad
              \sigma_r = \begin{cases}
                  \mathrm{RMSE}_r & \text{if available} \\
                  \sigma_r^{(0)}  & \text{otherwise}
              \end{cases}
          \end{equation}
          This allows the optimizer to weight edges according to the
          registration quality: well-aligned point clouds yield low RMSE
          and therefore small sigmas (high confidence), while poor
          alignments produce large sigmas (low confidence).
    \item \textbf{Inlier scaling (optional).}
          When enabled, both sigmas are scaled by the inverse square root
          of the inlier count~$n$:
          \begin{equation}
              \sigma_t \leftarrow \frac{\sigma_t}{\sqrt{n}}, \qquad
              \sigma_r \leftarrow \frac{\sigma_r}{\sqrt{n}}
          \end{equation}
          The motivation follows the standard error of the mean: given~$n$
          inlier correspondences, the uncertainty of the aggregate
          registration should decrease as $1/\sqrt{n}$.
          However, this scaling is \emph{disabled by default} because the
          point-level inliers from feature matching are not independent
          pose measurements---they share spatial correlations due to local
          surface geometry---and applying the scaling would
          over-weight high-overlap edges.
    \item \textbf{Clamping.}
          The resulting sigmas are clamped to configurable bounds to
          prevent numerically degenerate noise models:
          \begin{equation}
              \sigma_t \leftarrow
                  \mathrm{clamp}\!\left(\sigma_t,\;
                  \sigma_t^{\min},\; \sigma_t^{\max}\right), \qquad
              \sigma_r \leftarrow
                  \mathrm{clamp}\!\left(\sigma_r,\;
                  \sigma_r^{\min},\; \sigma_r^{\max}\right)
          \end{equation}
          with defaults
          $\sigma_t^{\min} = \SI{0.10}{\metre}$,
          $\sigma_t^{\max} = \SI{5.0}{\metre}$,
          $\sigma_r^{\min} = \SI{0.10}{\radian}$,
          $\sigma_r^{\max} = \pi\;\mathrm{rad}$.
          The lower bound ensures that even high-quality registrations do
          not produce artificially tight constraints, while the upper bound
          prevents effectively unconstrained edges from distorting the
          optimization.
\end{enumerate}

\paragraph{Robust noise model.}
After computing $\sigma_t$ and $\sigma_r$, a GTSAM diagonal noise model is
constructed.
By default, this base model is wrapped in a Huber
M-estimator~\cite{huber_robust} with threshold $k = 1.345$ to robustify
the optimization against outlier registrations.
The Huber loss function transitions from quadratic to linear penalization at
the threshold:
\begin{equation}
    \rho(r) =
    \begin{cases}
        \frac{1}{2} r^2                & \text{if } |r| \leq k \\
        k \left(|r| - \frac{k}{2}\right) & \text{if } |r| > k
    \end{cases}
\end{equation}
The value $k = 1.345$ is a standard choice that retains $95\%$ asymptotic
efficiency under a Gaussian distribution while bounding the influence of
outliers.
In the multi-robot calibration context, this means that a single
incorrect registration (e.g.\ from FRICP converging to a local minimum)
cannot dominate the cost function and corrupt the poses of correctly
calibrated robots.

\subsection{Optimization}

The resulting nonlinear least-squares problem is solved using the
Levenberg--Marquardt algorithm with a maximum of 100 iterations.
The optimized poses $\{\mathbf{X}_i^*\}$ are extracted and converted to
ROS\,2  \
\texttt{geometry\_msgs/TransformStamped} messages, which are
broadcast via the static TF broadcaster.

%% ---------------------------------------------------------------------------
\section{ROS\,2 Interfaces}
\label{sec:ros2_interfaces}

The MRC package defines three custom interface types to support
the distributed calibration workflow.

\subsection{KissFeatures Message}

The \texttt{mrc/msg/KissFeatures} message encapsulates the extracted
geometric features for a single robot:

\begin{itemize}
    \item \texttt{keypoints} (\texttt{PointCloud2}): 3D coordinates of the
          extracted keypoints.
    \item \texttt{descriptors} (\texttt{float32[]}): Flattened row-major array
          of FasterPFH descriptors with
          $|\texttt{descriptors}| = \texttt{keypoint\_count} \times \texttt{descriptor\_dim}$.
    \item \texttt{descriptor\_dim} (\texttt{uint32}): Dimensionality of each
          descriptor (33 for FasterPFH).
\end{itemize}

\subsection{GetRegistrationData Service}

The \texttt{mrc/srv/GetRegistrationData} service allows the master to request
preprocessed data from a slave.
The request specifies the requester's robot identifier and the selected
algorithm.
Depending on the algorithm, the response contains either a raw
\texttt{PointCloud2} (for FRICP) or a \texttt{KissFeatures} message
(for KISS-Matcher), along with metadata (robot ID, frame ID, timestamp).

\subsection{Calibrate Action}

The \texttt{mrc/action/Calibrate} action provides the primary interface for
triggering calibration with progress monitoring:

\begin{description}
    \item[Goal:] Target robot IDs (empty for all discovered robots),
                 pose graph usage flag, and topology selection.
    \item[Feedback:] Current calibration phase
                     (initializing, fetching data, registering, optimizing, publishing),
                     progress percentage, and a human-readable status message.
    \item[Result:] Overall success flag, total execution time, number of
                   robots calibrated and failed, and per-robot error metrics.
\end{description}

%% ---------------------------------------------------------------------------
\section{Summary}

The MRC system comprises approximately 3\,100 lines of C++ code distributed
across seven core source files plus header files and interface definitions.
\autoref{tab:source_files} provides an overview of the implementation
components and their responsibilities.

\begin{table}[htbp]
    \centering
    \caption{Overview of MRC source files and their responsibilities.}
    \label{tab:source_files}
    \begin{tabular}{l r l}
        \toprule
        \textbf{Source File} & \textbf{Lines} & \textbf{Responsibility} \\
        \midrule
        \texttt{mrc\_master\_node.cpp}         & 1\,336 & Orchestration, action server, TF publishing \\
        \texttt{mrc\_slave\_node.cpp}           &    510 & Data acquisition, feature extraction \\
        \texttt{kiss\_matcher\_backend.cpp}     &    299 & Feature-based registration \\
        \texttt{fricp\_backend.cpp}             &    214 & Dense ICP-based registration \\
        \texttt{icp\_refiner.cpp}               &    169 & Optional ICP refinement \\
        \texttt{registration\_factory.cpp}      &     24 & Backend instantiation (factory pattern) \\
        \texttt{pose\_graph\_optimizer.cpp}     &    187 & GTSAM-based graph optimization \\
        \bottomrule
    \end{tabular}
\end{table}

The modular design with interchangeable backends and optional optimization
stages allows systematic evaluation of different algorithmic configurations,
which is presented in the following chapter.

%% ---------------------------------------------------------------------------
%% Bibliography entries referenced in this chapter (add to your main .bib file):
%%
%% @article{kiss_matcher,
%%   title   = {{KISS-Matcher}: Fast and Robust Point Cloud Registration Revisited},
%%   author  = {Oh, Hyungtae and Lim, Hyoung-Kyu and Lim, Hyun and Myung, Hyun},
%%   journal = {IEEE Robotics and Automation Letters},
%%   year    = {2024}
%% }
%%
%% @inproceedings{fricp,
%%   title     = {Fast and Robust Iterative Closest Point},
%%   author    = {Zhang, Juyong and Yao, Yuxin and Deng, Bailin},
%%   booktitle = {IEEE Transactions on Pattern Analysis and Machine Intelligence},
%%   year      = {2022}
%% }
%%
%% @techreport{gtsam,
%%   title       = {{GTSAM}: General Smoothing and Mapping Library},
%%   author      = {Dellaert, Frank and {GTSAM Contributors}},
%%   institution = {Georgia Institute of Technology},
%%   year        = {2022},
%%   note        = {\url{https://gtsam.org}}
%% }
%%
%% @article{robin,
%%   title   = {{ROBIN}: a Graph-Theoretic Approach to Reject Outliers in Robust Estimation using Invariants},
%%   author  = {Shi, Jingnan and Yang, Heng and Carlone, Luca},
%%   journal = {arXiv preprint arXiv:2011.03659},
%%   year    = {2021}
%% }
%%
%% @book{huber_robust,
%%   title     = {Robust Statistics},
%%   author    = {Huber, Peter J. and Ronchetti, Elvezio M.},
%%   publisher = {John Wiley \& Sons},
%%   year      = {2009}
%% }
















